\subsection{Обобщение направления спуска}
\label{subsec:descent-direction}

\begin{remark}
Вектор $p$ является направлением спуска в точке $x_k$ тогда и только тогда,
когда
\[
  \nabla\T f(x_k)\, p < 0 ,
\]
и, в частности, годится всякий вектор вида
\[
  p_k = -F(x_k)\, \nabla f(x_k) ,
\]
где $F(x_k)$ — некоторая матрица. Это и есть идея методов второго порядка.
\end{remark}

\begin{idea}[title={Откуда берётся вид $p_k = -F \nabla f$}]
Подставим $p_k = -F \nabla f(x_k)$ в критерий спуска (рис.~\ref{fig:descent-cone}):
\[
  \nabla\T f(x_k)\, p_k = -\,\nabla\T f(x_k)\, F\, \nabla f(x_k) < 0
\]
для любой положительно определённой $F$ и любого $\nabla f(x_k) \ne 0$.
То есть \emph{любая} положительно определённая
матрица $F$ порождает корректное направление спуска, и выбор $F$ — это и есть
выбор метода:
\begin{itemize}[nosep,left=0pt]
  \item $F = E$ — обычный градиентный спуск;
  \item $F = H^{-1}(x_k)$ — метод Ньютона (\S\ref{subsec:second-order});
  \item $F \approx H^{-1}(x_k)$, построенная по накопленным градиентам, —
        квазиньютоновские методы.
\end{itemize}
Смысл множителя $F$ — исправить перекос линий уровня: антиградиент указывает
направление наибыстрейшего убывания в \emph{евклидовой} метрике, а $F$ меняет
метрику так, чтобы линии уровня стали ближе к сферам. Именно отношение $m/M$,
портившее все оценки выше, при удачном $F$ перестаёт играть роль.
\end{idea}

\subsection{Метод циклического покоординатного спуска}

Задача $f(x_1, x_2, \dots, x_n) \to \min$, начальное приближение
$\big(x_1^{(0)}, x_2^{(0)}, \dots, x_n^{(0)}\big)\T$.

Упрощаем задачу, минимизируя по одной координате:
\[
  f\big(x_1,\, x_2^{(0)},\, \dots,\, x_n^{(0)}\big) \to \min_{x_1} ,
  \qquad \exists\, \bar{x}_1 \ \text{— решение} .
\]
Далее либо
\[
  \left[
  \begin{aligned}
    &f\big(x_1^{(0)},\, x_2,\, \dots,\, x_n^{(0)}\big) \to \min_{x_2} , \\
    &f\big(\bar{x}_1,\, x_2,\, \dots,\, x_n^{(0)}\big) \to \min_{x_2} ,
  \end{aligned}
  \right.
\]
и так далее по всем координатам. Процедура зацикливается тогда, когда
\[
  \norm{x^{(k)} - x^{(k-1)}} < \delta .
\]

\begin{idea}
Две записанные альтернативы — это два разных режима:
\begin{itemize}[nosep,left=0pt]
  \item минимизация по $x_2$ при \emph{старом} значении $x_1^{(0)}$ — все
        одномерные задачи решаются независимо от одной и той же точки
        (параллельный вариант, аналог метода Якоби);
  \item минимизация при уже обновлённом $\bar{x}_1$ — найденные координаты сразу
        идут в дело (последовательный вариант, аналог метода Гаусса–Зейделя).
\end{itemize}
Второй вариант обычно сходится быстрее, и именно он подразумевается под
«циклическим» покоординатным спуском. Производные здесь не нужны вовсе — метод
нулевого порядка, а каждая одномерная подзадача решается любым методом
из~\S\ref{subsec:direct-search-1d}.
\end{idea}

\subsection{Метод Хука--Дживса}

Состоит из двух шагов.

\paragraph{Шаг 1. Исследующий поиск.}
Вычисляется значение $f$ в точке $x^{(0)}$; поочерёдно каждая переменная
изменяется на свой шаг из набора $\{h_1, \dots, h_n\}$.

Для первой координаты:
\[
  \begin{aligned}
    &f\big(x^{(0)} + h_1 e_1\big) < f\big(x^{(0)}\big)
     &&\Longrightarrow\quad x^{(0)} := x^{(0)} + h_1 e_1 ; \\[4pt]
    &f\big(x^{(0)} + h_1 e_1\big) \ge f\big(x^{(0)}\big)
     &&\Longrightarrow\quad
       \text{проверяем } f\big(x^{(0)} - h_1 e_1\big) < f\big(x^{(0)}\big) :
  \end{aligned}
\]
\[
  \begin{aligned}
    &\text{да} &&\Longrightarrow\quad x^{(0)} := x^{(0)} - h_1 e_1 ; \\
    &\text{нет} &&\Longrightarrow\quad \text{координата остаётся прежней.}
  \end{aligned}
\]
Затем то же самое для $x_2$, и так далее, пока не переберём все переменные.
Полученная точка $x^{(1)}$ — новая \textbf{базовая точка}.

\begin{idea}
Логика пробного шага по координате: сначала пробуем сдвиг «вперёд» $+h_i e_i$;
если значение не улучшилось — пробуем «назад» $-h_i e_i$; если и это не помогло —
координату не трогаем и переходим к следующей. За один исследующий поиск каждая
координата меняется не более одного раза, и значение $f$ только невозрастает.
Если после полного прохода базовая точка не сдвинулась вовсе, шаги $h_i$
уменьшают (например, вдвое) и повторяют поиск — это и есть критерий остановки.
\end{idea}

\paragraph{Шаг 2. Поиск по образцу.}
\[
  x^{(k+1)} = x^{(k)} + \big(x^{(k)} - x^{(k-1)}\big) ,
  \qquad\text{то есть}\qquad
  x^{(k+1)} = 2x^{(k)} - x^{(k-1)} .
\]
\[
  \begin{aligned}
    &f\big(x^{(k+1)}\big) < f\big(x^{(k)}\big)
     &&\Longrightarrow\quad \text{строим дальше, продолжаем поиск;} \\
    &f\big(x^{(k+1)}\big) > f\big(x^{(k)}\big)
     &&\Longrightarrow\quad
       \text{откатываемся, последняя удачная точка — } x^{(k)} .
  \end{aligned}
\]

\inputfig{figures/fig-hooke-jeeves}

\begin{idea}
Смысл второго шага: исследующий поиск за цикл выявил удачное направление
$x^{(k)} - x^{(k-1)}$ — значит есть смысл пройти по нему ещё раз, не тратя
вычислений на пробные шаги. Это дешёвая замена одномерной минимизации: вместо
подбора длины шага метод просто повторяет уже сделанное смещение, пока функция
убывает. Отсюда характерная траектория — короткие «ощупывающие» зигзаги,
перемежающиеся длинными прыжками по образцу.
\end{idea}
